\documentclass[11pt]{article}
% macro.tex loads enumitem and geometry twice, the second time with options.
% Loading them here first (with those options) avoids the fatal "Option clash".
\usepackage[shortlabels]{enumitem}
\usepackage[hmargin=2cm,vmargin=2cm]{geometry}

%__________________________________
%————|Pacakages|———————————————————
%----------------------------------
\usepackage[utf8]{inputenc}
\usepackage[T1]{fontenc}
\usepackage{bm}
\newcommand{\vect}[1]{\bm{#1}}
\usepackage{amsmath}
\usepackage[hmargin=2cm,vmargin=2cm]{geometry}
\usepackage{hyperref}
\usepackage{mathtools}
\usepackage{amssymb}
% enumitem is loaded below with [shortlabels] -- loading it here without
% options first would raise "Option clash for package enumitem".
\usepackage{authblk}
\newtheorem{property}{Property}
\usepackage{graphicx}
\usepackage{subcaption}
\usepackage{float}
\geometry{margin=1.5cm}
\usepackage{amsmath,amssymb,amsthm,mathtools}
\usepackage{mathrsfs,bm,color}
\usepackage{fancyhdr,lastpage}
\usepackage[shortlabels]{enumitem}
\usepackage{calc}
\usepackage{bbm}
\usepackage{graphicx}
\usepackage{upgreek}
\usepackage{lmodern}
\usepackage[normalem]{ulem}
\usepackage{verbatim}
\usepackage{bbm}
%\usepackage{ntheorem}
\usepackage{stmaryrd}
\usepackage{amsmath}
% \usepackage{amstex}
\usepackage{amssymb}
\usepackage{dsfont}
\usepackage{amsthm}
\usepackage{hyperref}
\hypersetup{
	colorlinks,
	linkcolor={red!50!black},
	citecolor={red!80!black},
	urlcolor={blue!80!black}
}
\usepackage[numbers]{natbib}
\usepackage{xcolor}

\usepackage{graphicx}

\usepackage{amsfonts}%
\usepackage{dsfont}

%--------------------------------
\def\ra{{\rightarrow}}
\newcommand{\Var}{\operatorname{Var}}


\parindent=1em
\allowdisplaybreaks
% geometry is loaded once at the top of this file (with these same options);
% \geometry{margin=1.5cm} above is what actually sets the margins.


%--Enumerate----


%___________________________________
%—|Eqs,Theorems,Lemma,Proofs, etc|——
%-----------------------------------

\numberwithin{equation}{section}

\usepackage{thmtools}
\declaretheorem[thmbox=M,name=Theorem,numberwithin=section]{theo}
\declaretheorem[name=Proposition,thmbox=M,numberwithin=section]{prop}
\declaretheorem[name=Lemma,thmbox=S,numberwithin=section]{lem}
%\declaretheorem[name=Definition,thmbox=M,numberlike=theorem]{defn}
\declaretheorem[name=Corollary,thmbox=M,numberwithin=section]{cor}
\declaretheorem[name=Conjecture,thmbox=M,numberwithin=section]{con}
\declaretheorem[name=Hypothesis,thmbox=S,numberwithin=section]{hyp}
\theoremstyle{definition}
\newtheorem{exam}{Example}[section]
\newtheorem{defn}{Definition}[section]
\newtheorem{rem}{Remark}[section]
\newtheorem{exo}{Exercise}[section]
\newtheorem{ass}{Assumption}[section]
\newtheorem{theorem}{Theorem}[section]
\newtheorem{lemma}[theorem]{Lemma}
\newtheorem{proposition}[theorem]{Proposition}
\newtheorem{corollary}[theorem]{Corollary}
\theoremstyle{plain}

\makeatletter
\renewenvironment{proof}[1][\proofname]{\par
	\pushQED{\qed}%
	\normalfont \topsep6\p@\@plus6\p@\relax
	\trivlist
	\item[\hskip\labelsep
	\sffamily\bfseries
	#1\@addpunct{.}]\ignorespaces
}{%
	\popQED\endtrivlist\@endpefalse
}
\makeatother


%__________________________________
%————|Commands/defs|———————————————
%----------------------------------

%--font--------------------
\newcommand{\jump}{{\vskip 0.3cm \noindent }}


%--english math mode-------
\newcommand{\for}{\mbox{ for }}
\newcommand{\with}{\mbox{ with }}
\renewcommand{\and}{\mbox{ and }}

%--usual functions---------
\newcommand{\e}{\mathrm{e}}

%--Field/usual sets--------
\newcommand{\N}{\mathbb{N}}
\newcommand{\Z}{\mathbb{Z}}
\newcommand{\Q}{\mathbb{Q}}
\newcommand{\R}{\mathbb{R}}
\newcommand{\C}{\mathbb{C}}
\renewcommand{\H}{\mathbb{H}}

\newcommand{\D}{\mathbb{D}}
\newcommand{\T}{\mathbb{T}}
\newcommand{\B}{\mathbb{B}}
\renewcommand{\S}{\mathbb{S}}
\newcommand{\Sn}{\mathbb{S}^{N-1}}

%--Analysis----------------
\newcommand{\dd}{\mathrm d}
\newcommand{\I}{\mathbb{I}}
%\newcommand{\Ind}[1]{\ensuremath{ \I_{\{#1\}}}}
\newcommand{\Ind}[1]{\1_{(#1)}}

%--Ortho. Pol--------------
\newcommand{\polclass}[1]{\ensuremath{\mathrm{#1}}}
\newcommand{\He}{\polclass{He}}
\newcommand{\Hek}{\polclass{He}_k}
\newcommand{\La}{\polclass{La}}
\newcommand{\Lak}{\polclass{La}_k}
\newcommand{\Ja}{\polclass{Ja}}
\newcommand{\Jak}{\polclass{Ja}_k}
\newcommand{\op}{\polclass{P}}
\newcommand{\opk}{\polclass{P}_k}


%--Proba-------------------
\newcommand{\E}{\mathbb{E}}
\renewcommand{\P}{\mathbb{P}}
\def\as{{ \mathrm{a.s.}  }}
\def\ed{\stackrel{{\mathcal{D}}}{=}}
\def\iid{\stackrel{{\mathrm{i.i.d}}}{\sim}}

%--Lin. Algebra------------

    %usual set-------
    \newcommand{\setsqrmat}[1]{\ensuremath{
     \mathrm{M}_{#1}(\R)
    }}
    \newcommand{\setsymmat}[1]{\ensuremath{
     \mathrm{M}^{(\mathrm{sym})}_{#1}(\R)
    }}
   
    
\newcommand{\<}{\ensuremath{ \langle }}
\renewcommand{\>}{\ensuremath{ \rangle }}
% \vect is already defined at the top of this file as \bm{#1} (bold italic).
% Switch to the upright-bold version below if you prefer it to match \mat/\tens,
% but then comment out the definition at the top instead of this one.
%\newcommand{\vect}[1]{\ensuremath{\boldsymbol{\mathbf{#1}}}}
\newcommand{\rdmvect}[1]{\ensuremath{\bm{#1}}}
\newcommand{\mat}[1]{\ensuremath{\boldsymbol{\mathbf{#1}}}}
\newcommand{\rdmmat}[1]{\ensuremath{\bm{#1}}}
\newcommand{\tens}[1]{ {\ensuremath{\underline{\boldsymbol{\mathbf{#1}}}}}}
\newcommand{\rdmtens}[1]{ {\ensuremath{\underline{\bm{#1}}}}}
\newcommand{\diag}[1]{\ensuremath{\mat{D}_{\vect{#1}} }}
\newcommand{\Tr}{\ensuremath{\mathrm{Tr}}}
\renewcommand{\det}{\ensuremath{\mathrm{det}}}
\newcommand{\rank}{\ensuremath{\mathrm{rank}}}
%•Norms
\newcommand{\norm}[1]{\ensuremath{ \|#1 \|}}
\newcommand{\normop}[1]{\ensuremath{ \|#1 \|_{\mathrm{op}} }} 

%\newcommand{\norm2}[1]{\ensuremath{ \|#1 \|_{2} }}
\newcommand{\normp}[1]{\ensuremath{ \|#1 \|_{p} }}
\newcommand{\normf}[1]{\ensuremath{ \|#1 \|_{\mathrm{F}} }}
\newcommand{\scalar}[1]{\ensuremath{ \< #1 \>}}

%--classical ensembles of RMT--
%•Gaussian ensembles
\newcommand{\classens}[1]{\ensuremath{\mathsf{#1}}}
\newcommand{\GE}{\ensuremath{{ \classens{GU_{\beta}E}}}}
\newcommand{\GOE}{\ensuremath{{ \classens{GOE}}}}
\newcommand{\GUE}{\ensuremath{{ \classens{GUE}}}}
\newcommand{\GSE}{\ensuremath{{ \classens{GSE}}}}
%•Laguerre ensembles
\newcommand{\LE}{\ensuremath{{ \classens{LU_{\beta}E}}}}
\newcommand{\LOE}{\ensuremath{{ \classens{LOE}}}}
\newcommand{\LUE}{\ensuremath{{ \classens{LUE}}}}
\newcommand{\LSE}{\ensuremath{{ \classens{LSE}}}}
%•Laguerre ensembles
\newcommand{\JE}{\ensuremath{{ \classens{JU_{\beta}E}}}}
\newcommand{\JOE}{\ensuremath{{ \classens{JOE}}}}
\newcommand{\JUE}{\ensuremath{{ \classens{JUE}}}}
\newcommand{\JSE}{\ensuremath{{ \classens{JSE}}}}
%•Circular ensembles
\newcommand{\CE}{\ensuremath{{ \classens{CU_{\beta}E}}}}
\newcommand{\COE}{\ensuremath{{ \classens{COE}}}}
\newcommand{\CUE}{\ensuremath{{ \classens{CUE}}}}
\newcommand{\CSE}{\ensuremath{{ \classens{CSE}}}}
%•Haar Groups
\newcommand{\grpUNb}{\ensuremath{\classens{U}_{\beta}(N)}}
\newcommand{\grpUN}{\ensuremath{ \classens{U}(N)}}
\newcommand{\grpON}{\ensuremath{ \classens{O}(N)}}
\newcommand{\grpSpN}{\ensuremath{ \classens{Sp}(N)}}

%--free proba-------------------------
\newcommand{\rectbox}{{\mathpalette\scaledboxplus\relax}}
\newcommand{\scaledboxplus}[2]{\scalebox{2}[0.95]{$#1\normalboxplus$}}
\newcommand{\recplus}{\fboxsep0pt\fbox{\rule{0.5em}{0pt}\rule{0pt}{0.5ex}}}

%--Partitions/permutations-----------
\newcommand{\classpartition}[1]{\ensuremath{\vect{#1}}}
\newcommand{\classpartitionset}[1]{\ensuremath{\mathcal{}{#1}}}
\newcommand{\partlambda}{\ensuremath{ \classpartition{\uplambda} }}
\newcommand{\partmu}{\ensuremath{ \classpartition{\upmu} }}
\newcommand{\partnu}{\ensuremath{ \classpartition{\upnu} }}
\newcommand{\partpi}{\ensuremath{ \classpartition{\uppi} }}

\newcommand{\partofint}[1]{\ensuremath{ \classpartitionset{P}{#1}}}
\newcommand{\partofk}{\ensuremath{ \classpartitionset{P}(k)}}
\newcommand{\nonpartofint}[1]{\ensuremath{ \classpartitionset{NC}(#1)}}
\newcommand{\nonpartofk}{\ensuremath{ \classpartitionset{NC}(k)}}
\newcommand{\sumpovern}{\ensuremath{\sum_{\partlambda \vdash n }}}
\newcommand{\sumpoverk}{\ensuremath{\sum_{\partlambda \vdash k }}}

\newcommand{\classpermutation}[1]{\ensuremath{#1}}
\newcommand{\permsigma}{\ensuremath{ \classpermutation{\upsigma} }}
\newcommand{\permtau}{\ensuremath{\classpermutation{\uptau}}}

%--symm. polynomials------------------
\newcommand{\classsympol}[1]{\ensuremath{\mathrm{#1}}}

\newcommand{\jack}{\classsympol{J}}
\newcommand{\jackpl}{\classsympol{J}_{\partlambda}}
\newcommand{\jackplb}{\classsympol{J}_{\partlambda}^{(\beta/2)}}
\newcommand{\jackpmu}{\classsympol{J}_{\partmu}}
\newcommand{\jackpmub}{\classsympol{J}_{\partmu}^{(\beta/2)}}
\newcommand{\jackpnu}{\classsympol{J}_{\partnu}}
\newcommand{\jackpnub}{\classsympol{J}_{\partnu}^{(\beta/2)}}

\newcommand{\schur}{\classsympol{s}}
\newcommand{\schurpl}{\classsympol{s}_{\partlambda}}
\newcommand{\schurpmu}{\classsympol{s}_{\partmu}}
\newcommand{\schurpnu}{\classsympol{s}_{\partnu}}
\newcommand{\LRcoeff}{c_{\partlambda,\partmu}^{\partnu}}
\newcommand{\JLRcoeff}{c_{\partlambda,\partmu}^{\partnu}(\beta/2)}


%--article-specific------------
\def\LocLip{{\mathrm{Lip_{loc}}  }}
\def\ks{{k_{\star}}}
\def\Yf{{ \rdmmat{Y}^{(f)}  }}

\def\coeffk{{\vartheta_k(f)}}
\def\coeffks{{\vartheta_{\ks}(f)}}
\def\varfZ{{\vartheta_0(f^2) - \vartheta_0(f)^2 }}
\def\coeffkn{{\vartheta_k(f_t)}}

\def\xk{{\vect{x}^k}}
\def\xks{{\vect{x}^\ks}}
\newcommand{\pert}{\mathrm{pert}}
\def\fm{{f_M}}
\def\fn{{f_t}}
\def\fdm{{f_{\delta,M}}}
\def\fdmp{{f_{\delta,M}^{\pert}}}
\def\fnp{{f_{t}^{\pert}}}
\def\fdmpt{{\Tilde{f}_{\delta,M}^{\pert} }}
\def\fnpt{{\Tilde{f}_{t}^{\pert} }}
%\def\Ytf{{  \rdmmat{\tilde Y}^{(f_L)}  }}
\def\Ytf{{\rdmmat{Y}_{L}^{(f_L)} }}
\def\Ytfdm{{ \rdmmat{\tilde Y}^{(\fdm)}  }}
\def\Ytfdmp{{ \rdmmat{\tilde Y}^{(\fdmp)}  }}
%\def\Ytfnp{{ \rdmmat{\tilde Y}^{(\fnp)}  }}
\def\Ytfnp{{\rdmmat{Y}_{L}^{(\fnp)} }}
\def\Ytfdmpt{{ \rdmmat{\tilde Y}^{(\fdmpt)}  }}
\def\Yfm{{ \rdmmat{Y}^{(\fm)}  }}
\def\Yfdm{{ \rdmmat{Y}^{(\fdm)}  }}
\def\Yfdmp{{ \rdmmat{Y}^{(\fdmp)}  }}
\def\Yfdmpt{{ \rdmmat{Y}^{(\fdmpt)}  }}

\def\wkZ{{w^{(k)}_Z}}
\def\tC{\widetilde{\mat C}}
\def\tY{\widetilde{\rdmmat Y}}

\def\coeffkfdm{{\vartheta_k(\fdm)}}
\def\coeffksfdm{{\vartheta_{\ks}(\fdm)}}
\def\coeffkfdmp{{\vartheta_k(\fdmp)}}
\def\coeffksfdmp{{\vartheta_{\ks}(\fdmp)}}
\def\coeffkfnp{{\vartheta_k(\fnp)}}
\def\coeffksfnp{{\vartheta_{\ks}(\fnp)}}

\def\Lpdx{{L^p(\dd x)}}
\def\Ltwodx{{L^2(\dd x)}}
\def\Lfourdx{{L^4(\dd x)}}

\def\LpZ{{L^p(\mu_Z)}}
\def\LoneZ{{L^1(\mu_Z)}}
\def\LtwoZ{{L^2(\mu_Z)}}
\def\LfourZ{{L^4(\mu_Z)}}
\def\LinfZ{{L^{\infty}(\mu_Z)}}

%--commands crafted for two channels mode--
\newcommand{\bkt}[2]{B^{({#1}), t}_{#2}}
\newcommand{\tbt}[1]{\vect{b}^{t}_{#1}}
\newcommand{\bttt}[0]{\tilde{\vect{b}}^{t}}
\newcommand{\akt}[2]{A^{({#1}), t}_{#2}}
\newcommand{\tat}[1]{\mat{A}^{t}_{#1}}
\newcommand{\ts}[2]{\mat{S}_{#1, #2}}
\newcommand{\trn}[2]{\mat{R}_{#1, #2}}
\newcommand{\tsig}[2]{\mat{\Sigma}^{#1}_{#2}}
\newcommand{\xat}[3]{\hat{x}^{({#1}), #2}_{#3}}
\newcommand{\xtt}[2]{\vect{x}^{#1}_{#2}}
\newcommand{\xttt}[1]{\tilde{\vect{x}}^{#1}}


\def\tib{\vect{b}}
\def\tia{\mat{A}}
\def\tix{\vect{x}}
\def\xone{x^{(1)}}
\def\xtwo{x^{(2)}}
\def\oner{\tilde{1}_R}
\def\onel{\tilde{1}_L}
\def\ovN{{\frac{1}{N}}}
\def\ovsN{{\frac{1}{\sqrt{N}}}}
\def\dovN{{\frac{\mat{\Delta}^{-1}}{N}}}
\def\calij{\mathcal{I}-\mathcal{J}}
\def\calijs{\mathcal{I}-\mathcal{J^*}}
\def\kerij{\mathrm{Ker}(\mathcal{I}-\mathcal{J})}
\def\kerijs{\mathrm{Ker}(\mathcal{I}-\mathcal{J}^*)}

%--commands crafted for tensor mode--
\newcommand{\Tind}[2]{T_{#1 ... #2}}

%--maths symbol commands--
\newcommand\normvert[1]{\left\lVert{#1}\right\rVert}
\newcommand{\para}[1]{\left( {#1} \right)}
\newcommand{\brac}[1]{\left[ {#1} \right]}


%--other commands-----


\newcommand\given[1][]{\:#1\vert\:}
\DeclareMathOperator{\1}{\mathbbm{1}}

\newcommand{\ba}{{\beta}}
\newcommand{\bb}{{\beta_{SNR}}}
\newcommand{\bc}{{\beta_{S}}}
\newcommand{\cR}{\mathcal{R}}
\newcommand{\cL}{\mathcal{L}}
\newcommand{\cG}{\mathcal{G}}
\newcommand{\cB}{\mathcal{B}}
\newcommand{\cS}{\mathcal{S}}
\newcommand{\cH}{\mathcal{H}}
\newcommand{\sD}{\mathscr{D}}
\newcommand{\sO}{\mathcal{O}}
\newcommand{\so}{\mathcal{o}}
\newcommand{\cQ}{\mathcal{Q}}
\newcommand{\sS}{\mathscr{S}}
\newcommand{\sE}{\mathscr{E}}
\newcommand{\sM}{\mathcal{M}}
\newcommand{\sP}{\mathscr{P}}
\newcommand{\cM}{\mathcal{M}}
\newcommand{\sC}{\mathscr{C}}
\newcommand{\cC}{\mathcal{C}}
\newcommand{\sI}{\mathscr{S}}
\newcommand{\sL}{\mathscr{L}}
\newcommand{\sX}{\mathcal{X}}
\newcommand{\cI}{{\mathcal I}}
\newcommand{\bpi}{\bm{\pi}}

\newcommand{\ncal}{\mathcal{N}}
\renewcommand{\P}{\mathbb{P}}
\newcommand{\mmm}{\mathrel{}\mid\mathrel{}}

\newcommand{\acts}{\curvearrowright}
\newcommand{\dx}{\dmesure\!}
\newcommand{\deron}[2]{\frac{\partial #1}{\partial #2}}
\newcommand{\derond}[2]{\frac{\partial^2 #1}{\partial #2 ^2}}
\newcommand{\deronc}[3]{\frac{\partial^2 #1}{\partial #2 \partial #3}}
\newcommand{\esp}[2][]{\mathbb{E}_{#1}\!\left[ #2 \right]}
\newcommand{\espcond}[3][]{\mathbb{E}_{#1}\!\left[ #2 \rule{0pt}{4mm}\mvert #3 \right]}
\newcommand{\loi}[1]{\,\dmesure\! P_{#1}}
\newcommand{\grad}{\gradient}
\newcommand{\intg}{[\![}
\newcommand{\intd}{]\!]}
\newcommand{\mvert}{\mathrel{}\middle|\mathrel{}}
%\newcommand{\norm}[1]{\left\lvert #1 \right\rvert}
\newcommand{\Norm}[1]{\left\lVert #1 \right\rVert}
\newcommand{\prsc}[2]{\left\langle #1\,, #2 \right\rangle}
\newcommand{\rot}{\rotationnel}
\newcommand{\trans}{\leftidx{^\text{t}}}
\newcommand{\ubvs}{\tilde{\bm{\sigma}}}
\newcommand{\ubvp}{\tilde{\bm{\rho}}}
\newcommand{\us}{\tilde{\sigma}}
\newcommand{\bR}{\bm{R}}
\newcommand{\btR}{\tilde{\bm{R}}}
\newcommand{\bQ}{\bm{Q}}
\newcommand{\bM}{\bm{M}}
\newcommand{\bS}{\bm{\Sigma}}
\newcommand{\bA}{\bm{A}}
\newcommand{\bO}{\bm{O}}
\newcommand{\bL}{\bm{\Lambda}}
\newcommand{\bD}{\bm{\Delta}}
\newcommand{\bd}{\bm{D}}
\newcommand{\bP}{\bm{P}}
\newcommand{\bU}{\bm{U}}
\newcommand{\bC}{\bm{C}}
\newcommand{\be}{\bm{E}}
\newcommand{\bE}{\bm{\mathcal{E}}}
\newcommand{\bz}{\bm{z}}
\newcommand{\bN}{\bm{N}}
\newcommand{\bB}{\bm{B}}
\newcommand{\bT}{\bm{T}}
\newcommand{\bI}{\bm{I}}
\newcommand{\bxi}{\bm{\xi}}
\newcommand{\bx}{{\bm{x}}}
\newcommand{\bu}{{\bm{u}}}
\newcommand{\bt}{{\bm{t}}}

\newcommand\bg{{\bm{g}}}
\newcommand{\out}{\mathrm{out}}

\DeclareMathOperator{\bij}{Bij}
\DeclareMathOperator{\diam}{diam}
\DeclareMathOperator{\Div}{div}
\DeclareMathOperator{\dmesure}{d}
\DeclareMathOperator{\gradient}{grad}
\DeclareMathOperator{\Id}{Id}
\DeclareMathOperator{\im}{Im}
\DeclareMathOperator{\rk}{rg}
\DeclareMathOperator{\rotationnel}{rot}
\DeclareMathOperator{\tr}{Tr}


\renewcommand{\leq}{\leqslant}
\renewcommand{\geq}{\geqslant}
\renewcommand{\epsilon}{\varepsilon}

%\renewcommand{\FrenchLabelItem}{\textbullet}
\def\supp{{\mbox{supp}}}
\def\pP{{\mathbb P}}

\def\pQ{{\mathbb Q}}
\def\pG{{\mathbb G}}
\def\cA{{\mathcal A}}
\newcommand{\iii}{i_1,\dots,i_p}
\def\by{{\mathbf{y}}}
\def\bs{{\mathbf{s}}}
\newcommand{\tx}{\tilde{x}}
\def\tr{{\rm Tr}}
\def\R{{\mathbb R}}
\def\rank{{\rm rank}}
\def\T{{\mathbb T}}
\def\IND{{\mathbf 1}}

% Commande pour créer un chiffre entouré avec épaisseur de ligne ajustée
\newcommand{\customcircled}[1]{
  \tikz[baseline=(char.base)]{
    \node[shape=circle, draw, inner sep=1pt, line width=0.8pt] (char) {#1};
  }
}

\definecolor{FireBrick}{rgb}{0.698, 0.133, 0.133}




\theoremstyle{remark}
\newtheorem{remark}[theorem]{Remark}


\title{\bf Continual Learning for TP4}
\author{Celia Budelot, Taliesin Perez}
\date{September 2026, IdePHICS Lab, EPFL}

\begin{document}
\maketitle
\tableofcontents
\newpage

\section{Introduction}
\subsection{Motivation}
\subsection{Summary of the project - goal}

\section{Theory and convention used}
\begin{defn} \label{def:spiked-wigner}
Let $\lambda \in \mathbb{R}$ and let $\vect x \in \mathbb{R}^n$ be a signal vector satisfying \( \|\vect x\|^2 = 1\). Let $\rdmmat Z$ be a symmetric Gaussian Wigner matrix, that is,
\begin{align}
Z_{ij} = Z_{ji}, \qquad Z_{ij} \overset{\mathrm{i.i.d.}}{\sim} \mathcal N(0,\frac{1}{n})    \quad \text{for } 1\le i<j\le n,
\end{align}
and
\begin{align}
Z_{ii} \overset{\mathrm{i.i.d.}}{\sim} \mathcal N(0,\frac{1}{n}) \quad \text{for } 1\le i\le n.
\end{align}
\noindent
We define the spiked model as
\begin{equation}
    \rdmmat Y = \lambda \,\vect x \vect x^T + \rdmmat Z,
    \label{norm2}
\end{equation}
\end{defn}
For two correlated spikes, we define:
\begin{defn}[Multi-Task Spiked Wigner Model]\leavevmode\\
    Let $\lambda_1, \lambda_2 \in \R$ the signal-to-noise ratio per task and $\vect x_1, \vect x_2 \in \R^n$, with $\norm{\vect x_1} = \norm{\vect x_2} = 1$ the correlated signals per task from a multivariate Gaussian Prior:
    \begin{align}
        \forall i \in[n]\,,\; \para{x_{1,i}\;\; x_{2,i}}^T \overset{\mathrm{i.i.d.}}{\sim} \mathcal{N}(\vect 0, \mat C)\, ,\quad\mathrm{with}\quad \mat C := \begin{pmatrix}
            1 & \rho \\
            \rho & 1 
        \end{pmatrix}\,,\quad \mathrm{and}\quad \rho \in [0,1]\,.
    \end{align} 
    We define independent  Gaussian noises $\rdmmat Z_1$ and $\rdmmat Z_2$, namely:
    \begin{align}
        \forall (t,i,j) \in\{1,2\}\times[n]^2 \,, \; Z_{t,ij} = Z_{t,ji} \overset{\mathrm{i.i.d.}}{\sim} \mathcal{N}\bigg(0,\frac{1}{n}\bigg) .
    \end{align}
    We write the multi-spiked model as:
    \begin{align}
        \rdmmat Y_1 \, = \, \lambda_1 \vect x_1 \vect x_1^T + \rdmmat Z_1 \, , \\
        \rdmmat Y_2 \, = \, \lambda_2 \vect x_2 \vect x_2^T + \rdmmat Z_2 \, .
    \end{align}
    \noindent
    Our observation is then:
    \begin{equation}
    \rdmmat Y := \alpha \rdmmat Y_1 + \sqrt{1-\alpha^2} \cdot \rdmmat Y_2
    \label{eq:2spikes}
    \end{equation}
    \noindent
    with $ \alpha$ the importance factor, giving how important the task $i = 1,2$ is.
\end{defn}

\subsection{ BBP transition and Spectral method} 

The goal is to recover the hidden vector $\vect x$ from the observation $\rdmmat Y$. The possibility of recovering the signal is governed by the signal-to-noise ratio parameter $\lambda$. When $\lambda$ is sufficiently large, an eigenvalue corresponding to the signal separates from the bulk of the spectrum. This phenomenon is known as the \emph{BBP phase transition} and indicates that the signal becomes detectable through spectral methods.
\noindent
Before describing this transition, it is important to recall that the global distribution of the eigenvalues of a Wigner matrix admits a deterministic limit.

\begin{theo}[Semicircle Law]\label{thm:semi-circle}
Let $\lambda_1(Y),\dots,\lambda_n(Y)$ denote the eigenvalues of $\rdmmat Y$. The empirical spectral distribution
\begin{align}
\hat{\mu}_Y := \frac{1}{n}\sum_{i=1}^n \delta_{\lambda_i(Y)}
\end{align}
converges weakly almost surely as $n \to \infty$ to the semi-circle distribution with density
\begin{align}
    \rho_{sc}(x) := \frac{1}{2\pi}\sqrt{4-x^2}\,\mathbf{1}_{|x|\le 2}.
\end{align}
\end{theo}


\noindent
Furthermore, spectral analysis does not only yield eigenvalues distribution, or equivalently, it does not only give a certification about whether one can recover a signal or not. 
\noindent
While one can obtain the threshold by computing the top eigenvalue of $\rdmmat Y$, the top eigenvector associated corresponds to an estimator $\hat {\vect x}$ of $\vect x$. 
\noindent
To know if this is a good estimation, We define the overlap: 
\begin{equation}
    m_n = \lvert \langle \hat{\vect x} 
     , \vect x\rangle \rvert = \rvert\sum_{i=1}^n \hat{x}_i x_i \rvert \, .
    \label{eq:overlap_u}
\end{equation}
\noindent
One can think about the overlap as the alignment between the estimator and the signal. If the estimator is good, they should be aligned and the overlap should go to $1$. Otherwise, if the estimator is randomly chosen, it is most likely orthogonal to the signal in high dimension and the overlap should go to $0$.
\medskip 

\begin{theo}[BBP transition]\label{thm:bbp}
Let $\hat{\vect x}$ be the top eigenvector of $\rdmmat Y$ associated to the eigenvalue $\lambda_{\max}$. Additionally, $\hat{\vect x}$ is normalised so that $\|\hat{\vect x}\| = 1$, consistently with $\|\vect x\| = 1$. Then, as $n\to\infty$, the largest eigenvalue satisfies
\begin{align}
\lambda_{\max}(Y)\xrightarrow{\mathrm{a.s.}}
\begin{cases}
2 & \text{if } \lambda \le 1,\\[0.3em]
\lambda+\lambda^{-1} & \text{if } \lambda > 1,
\end{cases}
\end{align}
and the overlap \eqref{eq:overlap_u} between $\hat{\vect x}$ and the spike direction $\vect x$ satisfies
\begin{align}
    m_n
    %
    \xrightarrow{\mathrm{a.s.}}
    %
    m := \begin{cases}
        0 & \text{if } \lambda \le 1,\\[0.3em]
        \sqrt{1-\lambda^{-2}} & \text{if } \lambda > 1.
    \end{cases}
\end{align}
In particular, weak recovery by PCA is possible if and only if \(\lambda>1\).
\end{theo}

\section{First Observations}

\subsection{Eigenvalues of the signal matrix}

\paragraph{Reduction to a spiked Wigner model:}
We study what the behaviour of a two correlated spike should theoretically be, and compare that to simulations.\\
The observable splits into a noise part and a low rank signal part:
\begin{equation}
    \rdmmat Y = \underbrace{a\,\vect{x_1}\vect{x_1}^T + b\,\vect{x_2}\vect{x_2}^T}_{=:\ \rdmmat P}
              + \underbrace{\alpha \rdmmat Z_1 + \sqrt{1-\alpha^2}\,\rdmmat Z_2}_{=:\ \rdmmat W}
    \qquad
    a := \alpha\lambda_1, \quad b := \sqrt{1-\alpha^2}\,\lambda_2
    \label{eq:Y-split}
\end{equation}
Since $\rdmmat Z_1$ and $\rdmmat Z_2$ are independent with entries of variance $1/n$, the entries of $\rdmmat W$ are Gaussian with variance $(\alpha^2 + 1 - \alpha^2)/n = 1/n$: the noise $\rdmmat W$ is again a Wigner matrix. So \eqref{eq:Y-split} is a rank-two spiked Wigner model whose spike matrix $\rdmmat P$ carries all the dependence on $\alpha$, $\lambda_1$, $\lambda_2$ and $\rho$. We can then compute its spectrum in closed form.

\medskip
\noindent
\textbf{Which matrix are we working with?} The whole section alternates between two different matrices, and it is worth fixing the vocabulary once and for all:
\begin{itemize}[topsep=2pt,itemsep=1pt]
    \item the \emph{noiseless} (or signal) matrix $\rdmmat P$, of rank at most two. It is not observable, it is the object we can diagonalise exactly, in closed form, because it contains no randomness beyond the signals themselves;
    \item the \emph{noisy} (or observed) matrix $\rdmmat Y = \rdmmat P + \rdmmat W$. This is what we actually measure, and its spectrum can only be described asymptotically, through random matrix theory.
\end{itemize}
The strategy is: solve $\rdmmat P$ exactly, then feed that solution into the BBP theory to get $\rdmmat Y$. To keep the two apart, we reserve a distinct set of symbols for each:
\begin{center}
\begin{tabular}{l|l|l}
 & \textbf{noiseless} $\rdmmat P$ & \textbf{noisy} $\rdmmat Y = \rdmmat P + \rdmmat W$ \\
\hline
eigenvalues     & $\theta_\pm$ (exact, \eqref{eq:theta-pm})        & $\lambda^{\mathrm{out}}_\pm$ (asymptotic, \eqref{eq:bbp}) \\
eigenvectors    & $\hat{\vect v}_\pm$ (exact, \eqref{eq:v-final})  & $\hat{\vect u}_\pm$ (asymptotic, \eqref{eq:bbp-vec}) = $\hat{\vect x}$ for naïve method only \\
overlaps        & $|\langle \hat{\vect v}_\pm, \vect{x_i}\rangle|$, \eqref{eq:overlaps-signals} & $m_i = |\langle \hat{\vect x}, \vect{x_i}\rangle|$, \eqref{eq:m12} \\
\end{tabular}
\end{center}
\noindent
Note finally that $\lambda_1, \lambda_2$ always denote the signal-to-noise ratios of the two tasks, never eigenvalues; eigenvalues of $\rdmmat P$ are written $\theta$.

\paragraph{Factorised form:}
Since $\rdmmat P$ is built from two rank 1 signals, it has a rank of at most two and its spectrum is determined by the plane $\operatorname{span}(\vect{x_1},\vect{x_2})$. We write
\begin{equation}
    \rdmmat P = \rdmmat X \rdmmat K \rdmmat X^T
    \qquad
    \rdmmat X := \big( \vect{x_1} \ \ \vect{x_2} \big) \in \mathbb{R}^{n\times 2}
    \qquad
    \rdmmat K := \begin{pmatrix} a & 0 \\ 0 & b \end{pmatrix}
    \label{eq:P-factor}
\end{equation}
And using that:
\begin{equation}
    \det(\rdmmat{AB} + \rdmmat I) = \det(\rdmmat{BA} + \rdmmat I) 
\end{equation}
and applying is using that $\rdmmat A = \rdmmat X$ and $\rdmmat B = \rdmmat K\rdmmat X^T$, we have that:
\begin{equation}
    \det(\rdmmat P - \theta\, \rdmmat I_n) = (-\theta)^n \det\big(\rdmmat I_n - \rdmmat{AB}/\theta\big)
    = (-\theta)^n \det\big(\rdmmat I_2 - \rdmmat{BA}/\theta\big) ,
\end{equation}
so that, for $\theta \neq 0$, $\theta$ is an eigenvalue of the noiseless $\rdmmat P$ if and only if $\det(\rdmmat I_2 - \rdmmat{BA}/\theta) = 0$.
and as $\rdmmat{BA} = \rdmmat{KX^TX} $ so the non-zero eigenvalues of $\rdmmat P$ are exactly the eigenvalues of the $2\times 2$ matrix
\begin{equation}
    \rdmmat K \rdmmat G,
    \qquad
    \rdmmat G := \rdmmat X^T \rdmmat X
    = \begin{pmatrix} \|\vect{x_1}\|^2 & \langle \vect{x_1},\vect{x_2}\rangle \\ \langle \vect{x_1},\vect{x_2}\rangle & \|\vect{x_2}\|^2 \end{pmatrix}
    = \begin{pmatrix} 1 & \hat\rho \\ \hat\rho & 1 \end{pmatrix},
    \qquad
    \hat\rho := \langle \vect{x_1},\vect{x_2}\rangle
    \label{eq:gram}
\end{equation}
where $\rdmmat G$ is defined as shown above and we used $\|\vect{x_1}\| = \|\vect{x_2}\| = 1$. The remaining $n-2$ eigenvalues of $\rdmmat P$ are zero.

\paragraph{Empirical vs theoretical overlap:}
As the spikes are correlated, the pairs $(x_{1,i},x_{2,i})$ are i.i.d.\ with $\mathbb{E}[x_{1,i}x_{2,i}] = \rho$, so by the law of large numbers the empirical overlap concentrates around the correlation $\rho$:
\begin{equation}
    \hat\rho = \frac{\langle \vect{x_1},\vect{x_2}\rangle}{\|\vect{x_1}\|\,\|\vect{x_2}\|} \xrightarrow[n\to\infty]{\text{a.s.}} \rho
    \label{eq:overlap}
\end{equation}
In the high-dimensional limit, we may therefore replace $\rdmmat G$ by the covariance matrix $\rdmmat C$ defined in \eqref{eq:2spikes}.

\paragraph{The eigenvalues of the noiseless matrix $\rdmmat P$:}
These are the $\theta$'s, not to be confused with the eigenvalues of the observed $\rdmmat Y$ treated in \eqref{eq:bbp}. From \eqref{eq:gram}, we find that:
\begin{equation}
    \det(\rdmmat I_2 - \rdmmat K\rdmmat G/\theta)
    = \frac{1}{\theta^2}\det(\theta\,\rdmmat I_2 - \rdmmat K\rdmmat G)
    = \frac{1}{\theta^2}\det\begin{pmatrix} \theta - a & -a\rho \\ -b\rho & \theta - b \end{pmatrix}
    = 0
\end{equation}
can be solved using that:
\begin{equation}
    \rdmmat K \rdmmat G = \begin{pmatrix} a & a\rho \\ b\rho & b \end{pmatrix},
    \qquad
    \operatorname{tr}(\rdmmat K\rdmmat G) = a + b,
    \qquad
    \det(\rdmmat K\rdmmat G) = ab\,(1-\rho^2),
\end{equation}
so the characteristic polynomial is $\theta^2 - (a+b)\,\theta + ab(1-\rho^2) = 0$, whose roots are
\begin{equation}
    \boxed{\;
    \theta_\pm = \frac{a+b}{2} \pm \sqrt{\Big(\frac{a-b}{2}\Big)^2 + ab\,\rho^2}
    \;}
    \qquad
    a = \alpha\lambda_1, \quad b = \sqrt{1-\alpha^2}\,\lambda_2
    \label{eq:theta-pm}
\end{equation}
The discriminant is non-negative, so $\theta_\pm \in \mathbb{R}$. These two numbers are exact for every $n$: no noise and no large-$n$ limit has been used so far.\\
\textcolor{red}{We can then plot the empirical highest eigenvalues $\lambda_1$ and $\lambda_2$ with the theoretical ones $\theta_+$ and $\theta_-$ to compare them (see plots).}\\
\textcolor{purple}{TRUE FOR $\lambda_1 = \lambda_2=1$: We also note that when solving for $\theta_+ = 1$ and $\theta_- = 1$ we find $\lambda^*_1 = \frac{1}{1 + \rho} \frac{1}{\sqrt{\frac{1}{2}}}$ with the $\sqrt{\frac{1}{2}}$ coming from equally weighted spikes. We also find $\lambda^*_2 = \frac{1}{1 - \rho} \frac{1}{\sqrt{\frac{1}{2}}}$. We therefore note that those value of $\lambda^*_1$ and $\lambda^*_2$ will be the point at which the BBP transition will occur. We also note that in the limit $\rho = 0$, we find that $\lambda_1^*=\lambda^*_2=\sqrt{2}$, what does that mean ?}

\paragraph{The eigenvectors of the noiseless matrix $\rdmmat P$:}
We now derive the eigenvectors of $\rdmmat P$ in closed form. We assume $a, b > 0$ (i.e.\ $\alpha \in (0,1)$ and $\lambda_1, \lambda_2 > 0$) and $\rho \in [0,1)$, so that $\det(\rdmmat K \rdmmat G) = ab(1-\rho^2) > 0$ and both $\theta_\pm$ are strictly positive. We introduce
\begin{equation}
    \Delta := \sqrt{\Big(\frac{a-b}{2}\Big)^2 + ab\,\rho^2},
    \qquad \text{so that} \qquad
    \theta_\pm = \frac{a+b}{2} \pm \Delta,
    \label{eq:Delta}
\end{equation}
and we assume $\Delta > 0$, i.e.\ we exclude the degenerate case $a = b$, $\rho = 0$, in which $\theta_+ = \theta_-$ and every vector of $\operatorname{span}(\vect{x_1},\vect{x_2})$ is an eigenvector.

\paragraph{Going to the plane of the signals (méthode de Celia):} Let $\vect v$ be a unit eigenvector of $\rdmmat P$ with eigenvalue $\theta \neq 0$. Since
\begin{equation}
    \rdmmat P \vect v = a\,\langle \vect{x_1}, \vect v\rangle\,\vect{x_1} + b\,\langle \vect{x_2}, \vect v\rangle\,\vect{x_2} = \theta\,\vect v ,
\end{equation}
$\vect v$ must lie in $\operatorname{span}(\vect{x_1},\vect{x_2})$, and we write
\begin{equation}
    \vect v = c_1\,\vect{x_1} + c_2\,\vect{x_2},
    \qquad
    \langle \vect{x_1}, \vect v\rangle = c_1 + \rho\, c_2,
    \qquad
    \langle \vect{x_2}, \vect v\rangle = \rho\, c_1 + c_2 .
    \label{eq:ansatz}
\end{equation}
Substituting into the eigenvalue equation and identifying the coefficients of $\vect{x_1}$ and $\vect{x_2}$ (which are linearly independent since $\rho < 1$) gives
\begin{equation}
    \begin{cases}
        a\,(c_1 + \rho\, c_2) = \theta\, c_1 ,\\[2pt]
        b\,(\rho\, c_1 + c_2) = \theta\, c_2 ,
    \end{cases}
    \qquad \text{i.e.} \qquad
    (\rdmmat K \rdmmat G - \theta\, \rdmmat I_2)\begin{pmatrix} c_1 \\ c_2 \end{pmatrix} = \vect 0 .
    \label{eq:eig-system}
\end{equation}
A non-trivial solution exists if and only if $\theta \in \{\theta_+, \theta_-\}$, recovering \eqref{eq:theta-pm}. For $\rho > 0$ we have $(\theta_\pm - a)(\theta_\pm - b) = ab\rho^2 > 0$, so $\theta_\pm \neq a$, and the first row yields
\begin{equation}
    c_1 = r_\pm\, c_2,
    \qquad
    r_\pm := \frac{\rho\, a}{\theta_\pm - a} = \frac{\theta_\pm - b}{b\rho},
    \label{eq:ratio}
\end{equation}
where the two expressions for $r_\pm$ agree precisely because of the characteristic equation. For the same reason, the second row of \eqref{eq:eig-system} is then automatically satisfied.

\paragraph{Normalization:} Using $\norm{\vect{x_1}} = \norm{\vect{x_2}} = 1$ and $\langle \vect{x_1},\vect{x_2}\rangle = \rho$,
\begin{equation}
    \norm{\vect v}^2 = c_1^2 + 2\rho\, c_1 c_2 + c_2^2
    = c_2^2\,\big( r_\pm^2 + 2\rho\, r_\pm + 1 \big) = 1
    \quad \Longrightarrow \quad
    c_2 = \pm \big( r_\pm^2 + 2\rho\, r_\pm + 1 \big)^{-1/2} .
\end{equation}
The bracket is strictly positive, since $r^2 + 2\rho r + 1 = (r+\rho)^2 + 1 - \rho^2 > 0$.

\paragraph{Final expression for the eigenvectors of the noiseless $\rdmmat P$:} Combining \eqref{eq:ratio} with the normalization, the unit-norm eigenvectors of $\rdmmat P$ associated with $\theta_\pm$ are, up to an overall sign,
\begin{equation}
    \hat{\vect v}_\pm = c_1^\pm\,\vect{x_1} + c_2^\pm\,\vect{x_2},
    \qquad
    \big(c_1^\pm,\, c_2^\pm\big) = \frac{\big(r_\pm,\, 1\big)}{\sqrt{r_\pm^2 + 2\rho\, r_\pm + 1}},
    \qquad
    r_\pm = \frac{\rho\, a}{\theta_\pm - a} .
    \label{eq:v-final}
\end{equation}

\paragraph{Overlaps of the noiseless eigenvectors with the signals:} The overlaps of $\hat{\vect v}_\pm$ with the signals follow directly from \eqref{eq:ansatz} and \eqref{eq:eig-system}, since $c_1 + \rho c_2 = \theta c_1/a$ and $\rho c_1 + c_2 = \theta c_2/b$:
\begin{equation}
    \langle \hat{\vect v}_\pm, \vect{x_1}\rangle = \frac{\theta_\pm}{a}\, c_1^\pm,
    \qquad
    \langle \hat{\vect v}_\pm, \vect{x_2}\rangle = \frac{\theta_\pm}{b}\, c_2^\pm .
    \label{eq:overlaps-signals}
\end{equation}
Eq. \eqref{eq:overlaps-signals} shows that the \emph{coefficients} $c_i^\pm$ (how much of each signal we put into the mixture) are not the same thing as the \emph{overlaps} $\langle \hat{\vect v}_\pm, \vect{x_i}\rangle$ (how much of each signal we actually see), because $\vect{x_1}$ and $\vect{x_2}$ are not orthogonal.

\medskip
\noindent
Multiplying the first relation in Eq. \eqref{eq:overlaps-signals} by $a\,\langle \hat{\vect v}_\pm, \vect{x_1}\rangle$, its second by $b\,\langle \hat{\vect v}_\pm, \vect{x_2}\rangle$, adding, and using the normalization $c_1^2 + 2\rho c_1 c_2 + c_2^2 = 1$, we get
\begin{equation}
    \theta_\pm \;=\; \hat{\vect v}_\pm^{\,T}\, \rdmmat P\, \hat{\vect v}_\pm
    \;=\; a\,\langle \hat{\vect v}_\pm, \vect{x_1}\rangle^2 + b\,\langle \hat{\vect v}_\pm, \vect{x_2}\rangle^2 .
    \label{eq:theta-pm-relation}
\end{equation}

\paragraph{TO PROOVE:}
We now study the noisy matrix $\rdmmat Y = \rdmmat P + \rdmmat W$. As $\rdmmat P$ has rank at most two, it can only create at most two outliers. Also, $\rdmmat W$ is a Wigner matrix, hence orthogonally invariant, so no direction of $\R^n$ is singled out by the noise. Therefore, the outliers can depend on $\rdmmat P$ only through its non-zero eigenvalues $\theta_\pm$, and never through the directions $\hat{\vect v}_\pm$ \textcolor{red}{WHY ?} This is what allows us to reuse the rank-one theory of Theorem~\ref{thm:bbp}.


\begin{equation}
    \lambda_\pm^{\mathrm{out}} \xrightarrow[n\to\infty]{\text{a.s.}}
    \begin{cases}
        \theta_\pm + \dfrac{1}{\theta_\pm}, & \theta_\pm > 1, \\[6pt]
        2, & \theta_\pm \le 1 ,
    \end{cases}
    \label{eq:bbp}
\end{equation}
which is the BBP transition applied to each effective signal separately. So that the two correlated tasks behave like two \emph{independent} rank-one spiked Wigner models with \emph{strengh} $\theta_+$ and $\theta_-$, and \eqref{eq:bbp} is exactly Theorem~\ref{thm:bbp} with $\lambda$ replaced by $\theta_\pm$.

\paragraph{The eigenvectors of the noisy matrix $\rdmmat Y$:}
Let $\hat{\vect u}_\pm$ denote the unit eigenvector of $\rdmmat Y$ associated with $\lambda_\pm^{\mathrm{out}}$; it is the noisy counterpart of the noiseless $\hat{\vect v}_\pm$. \textcolor{red}{TO PROOVE:}
\begin{equation}
    |\langle \hat{\vect u}_\pm, \hat{\vect v}_\pm \rangle|^2 \xrightarrow[n\to\infty]{\text{a.s.}} 1 - \frac{1}{\theta_\pm^{2}},
    \qquad
    \langle \hat{\vect u}_\pm, \hat{\vect v}_\mp \rangle \xrightarrow[n\to\infty]{\text{a.s.}} 0 ,
    \label{eq:bbp-vec}
\end{equation}
valid for $\theta_\pm > 1$. For $\theta_\pm \le 1$ the overlap vanishes.

\paragraph{Overlap of the estimator with the signals.}
The spectral estimator of the signals is the top eigenvector $\hat{\vect x} := \hat{\vect u}_+$ of the noisy matrix. Its overlap with $\vect{x_1}$ and $\vect{x_2}$ now follows by combining the noiseless geometry \eqref{eq:overlaps-signals} with the noisy damping \eqref{eq:bbp-vec}. Since $\hat{\vect v}_+$ and $\hat{\vect v}_-$ form an orthonormal basis of $\operatorname{span}(\vect{x_1},\vect{x_2})$, each signal decomposes as
\[
    \vect{x_i} = \langle \vect{x_i}, \hat{\vect v}_+\rangle\, \hat{\vect v}_+ + \langle \vect{x_i}, \hat{\vect v}_-\rangle\, \hat{\vect v}_- ,
\]
and taking the scalar product with $\hat{\vect x}$, the component of $\hat{\vect x}$ orthogonal to the plane contributes nothing,
\[
    \langle \hat{\vect x}, \vect{x_i}\rangle
    =\langle \hat{\vect u}_+, \hat{\vect v}_+\rangle
      \langle \hat{\vect v}_+, \vect{x_i}\rangle
    \;+\; \langle \hat{\vect u}_+, \hat{\vect v}_-\rangle
      \langle \hat{\vect v}_-, \vect{x_i}\rangle .
\]
Seeing that $\langle \hat{\vect u}_+, \hat{\vect v}_+\rangle\to\ \sqrt{1-\theta_+^{-2}}\ \text{by }\eqref{eq:bbp-vec}$ and $\langle \hat{\vect u}_+, \hat{\vect v}_-\rangle \to\ 0\ \text{by }\eqref{eq:bbp-vec}$. Substituting $\langle \hat{\vect v}_+, \vect{x_1}\rangle = (\theta_+/a)\,c_1^+$ and $\langle \hat{\vect v}_+, \vect{x_2}\rangle = (\theta_+/b)\,c_2^+$ from \eqref{eq:overlaps-signals},
\begin{equation}
    m_1 := |\langle \hat{\vect x}, \vect{x_1}\rangle| \xrightarrow[n\to\infty]{\text{a.s.}} \sqrt{1-\theta_+^{-2}}\;\frac{\theta_+}{a}\,\lvert c_1^+ \rvert ,
    \qquad
    m_2 := |\langle \hat{\vect x}, \vect{x_2}\rangle| \xrightarrow[n\to\infty]{\text{a.s.}} \sqrt{1-\theta_+^{-2}}\;\frac{\theta_+}{b}\,\lvert c_2^+ \rvert ,
    \label{eq:m12}
\end{equation}
for $\theta_+ > 1$, and $m_1, m_2 \to 0$ for $\theta_+ \le 1$.

\end{document}